\section{Discussion}\subsection{What the method changes in an audit}
The auditor first states the contrast, required precision, available
information and retrieval options. A sufficient initial interval
needs no query. Otherwise, a bulk response or fixed batch can avoid feedback;
adaptive retrieval is useful when it meets the required guarantee with fewer
permitted readings.

\paragraph{What the application establishes.}
With two readings shared across six comparisons, feedback guarantees the
required precision on 165 original and 57 external blocks where every
randomized fixed batch has worst-case success at most one half. Matching
batch constructions attain that bound. This is a change in the assurance
available at the same budget. The worked example also reports the six
interval widths in original loss units (Table~\ref{tab:absolute_joint_widths}).
These precision certificates complement the constructed mean-risk theorem;
they do not measure its $1/2$-versus-$2/3$ gap on the data. Across the
complete panels, recorded gains over the best fixed worst-width
batch are 7.81 and 11.11 percentage points, with smaller gains of 2.20 and
2.15 points over greedy adaptation. We do not infer improved forecasting
accuracy or operating savings from these precision results.

The separate-pair replay has smaller recorded gains of 0.407 and
0.119 percentage points. Both analyses use three systems per source;
repeated days and model pairs do not establish population prevalence.

The capacity controls show that feedback's larger gains need not survive
additional hourly information, while internal-reading errors can increase
the sufficient budget. These are alternative information models, not
deployment-calibrated bounds (Table~\ref{tab:measurement-assumptions}).
The half-width target limits the remaining range of each loss contrast;
we have not validated it as a threshold for financial action or model selection.
A secondary shared-query check
certifies a best model on 123 original and 207 external blocks, versus
122 and 201 for fixed batches and 118 and 208 for greedy adaptation.
On the external panel, greedy certifies one more best model despite its
lower joint-precision success. The width objective therefore does not
uniformly improve model choice.
Requests until a decision are another objective; neither theorem optimizes it.

\paragraph{Choosing an acquisition interface.}
Bulk uses one request and is faster on the measured three-task API workload,
so the implementation selects it when admissible. Under the separate
restricted-reading contract, it uses 22 two-reading and seven three-reading
prescriptions. Their 65 reserved readings compare with 87 for fixed batches,
but they require 51 information rounds rather than 29. The declared cost
frontier explains when each prescription is preferable; it is not measured
financial or service-wide benefit. Capacity bounds can likewise make fixed
retrieval sufficient. The practical output is a checked choice under a stated
contract, including cases where selective adaptation is unnecessary.

\paragraph{Breadth of the theory and limits of the implementations.}
The scalar classification permits compact convex sets, affine targets, a
finite linear-query dictionary and answer-independent legal query sets.
It includes signed targets, overlapping constraints and unequal fixed costs.
Fixed bounded errors can enter the hidden record (Online Resource~1, S6).
The finite dictionary matters: unrestricted linear measurements would permit
querying the scalar target directly. The algorithms are more specialized:
general polyhedra with an enumerated batch family or two-query planning;
cumulative sums with a known nonnegative total; and compact coordinate
recovery with disjoint exact totals and unit costs. They do not provide an
efficient arbitrary-budget precision optimizer for every convex set.

\section{Conclusions}
The risk objective determines whether feedback can improve the worst-case
guarantee for a model-comparison interval. Concavity gives a complete
no-improvement classification; nonconcave penalties admit sharp query-based
separations. Simultaneous comparisons have a different boundary: the largest
normalized width can benefit from feedback even under a mean-risk objective.
An exact two-query rule and checkable certificates turn that
distinction into an acquisition decision. The reference implementation chooses among supplied bulk, fixed and adaptive
certificates using declared access and resource limits. New comparisons require
new certificates; a catalogue exclusion alone is not an impossibility proof.
The public-data results show both when feedback changes assurance and when
bulk or fixed retrieval is preferable. These choices remain conditional on
the declared uncertainty set and access rules.
